\section{Introduction}\label{sec:intro}

Foodborne outbreak surveillance in the United States is a primary national source on how frequently illnesses associated with specific enteric pathogens result in hospitalization. Since 2009, state, local and territorial health departments have reported outbreaks across transmission modes to the National Outbreak Reporting System (NORS)~\cite{hall2013acute}. Descriptive summaries routinely compare pathogens by reported hospitalization percentages: in enteric outbreaks during 2009--2019, reported hospitalization was 20.9\% for \textit{Salmonella} and 22.8\% for \textit{Escherichia coli} outbreaks versus 2\% for norovirus outbreaks~\cite{wikswo2022enteric}, and annual surveillance reports tabulate hospitalization percentages alongside etiology~\cite{cdc2026surveillance2019}. Outbreak hospitalization metrics also enter composite severity indices~\cite{simpson2022analysing} and inform national burden estimates, for some pathogens serving as the sole empirical source~\cite{scallanwalter2025foodborne}.

For each outbreak, NORS records the estimated number of ill persons, the number hospitalized and, separately, the number of ill persons with known hospitalization status~\cite{cdc2014norsguidance}. The difference between total illnesses and persons with known status represents cases with unrecorded hospitalization outcomes rather than confirmed non-hospitalized cases. Published analyses resolve these unrecorded cases by convention: dividing hospitalizations by all ill persons implicitly treats every unrecorded case as not hospitalized~\cite{heiman2015escherichia,cdc2026surveillance2019}, whereas dividing by persons with known status assumes that unrecorded cases were hospitalized at the same rate as recorded cases~\cite{burke2019norovirus}. Although prior audits have examined field-level completeness in NORS~\cite{zhang2022evaluating} and demonstrated that structural missingness can bias derived outbreak metrics~\cite{diemer2024missingness}, neither the impact of unrecorded outcomes on pairwise pathogen rankings nor the joint statistical inference required when outbreak clustering, estimated standardization weights and multiple testing are present has been quantified.

Comparisons across pathogens face two further methodological challenges. First, pathogen groups occur in distinct food-preparation settings that independently predict severe outcomes~\cite{burke2019norovirus}, while state reporting intensity varies substantially~\cite{white2022foodborne} and shifted during the COVID-19 pandemic~\cite{ray2021decreased}; crude hospitalization percentages therefore conflate pathogen differences with setting and period mix. Second, cases cluster strongly within outbreaks whose sizes are heavily right-skewed: ignoring outbreak clustering or treating sample worst-case intervals as deterministic ignores sampling variability, whereas standardizing to a pooled reference distribution introduces dependence across groups through shared estimated weights.

To address these gaps simultaneously, we developed an integrated partial-identification and stacked M-estimation inference framework~\cite{manski1989anatomy,horowitz2000nonparametric,keiding2014standardization,ratitch2013missing,stefanski2002calculus,imbens2004confidence,romano2005exact} for six major foodborne pathogen groups in NORS during 2009--2023: \textit{Salmonella}, Shiga toxin--producing \textit{E. coli} (STEC), \textit{Campylobacter}, \textit{Clostridium perfringens}, norovirus genogroup I (GI) and norovirus genogroup II (GII). Beyond characterizing sample worst-case bounds and tipping points, we derive outbreak-cluster sandwich covariances that couple reference-weight uncertainty with within-cell sampling variance, construct Imbens--Manski uniformly valid 95\% confidence intervals and stepdown family-wise error rate (FWER) adjusted confidence tipping points across all 15 pairwise comparisons, diagnose outbreak-size concentration, and propose a Value-of-Information (VoI) follow-up allocation rule validated by retrospective masking on complete-status NORS outbreaks.
